\section{Prisma a facce parallele}
Quando un raggio luminoso attraversa un prisma a facce parallele, subisce due rifrazioni come in Figura 3.
\begin{figure}[H]
    \centering
    \begin{tikzpicture}[scale=1.25]
        % Definizione delle coordinate (A e B sono la base inferiore)
        \coordinate (A) at (0,0);       % Angolo retto inferiore sinistro
        \coordinate (B) at (0,5);       % Angolo retto inferiore destro
        \coordinate (C) at (2.5,2.5);       % Angolo superiore destro (inizio lato obliquo)
        \coordinate (D) at (2.5,0); 
        \coordinate (E) at (-1.5, 1.7);
        \coordinate (F) at (0,1.7);
        \coordinate (G) at (-1.5, 2.3);
        \coordinate (H) at (4,1.2);
        \coordinate (I) at (2.5, 1.2);
        \coordinate (J) at (4, 0.6);
        % --- DISEGNO DEL TRAPEZIO RETTANGOLO ---
        \draw[thick, fill=blue!5] (A) -- (B) -- (C) -- (D) -- cycle;
        \pic [draw, angle radius=1cm, angle eccentricity=1.6, fill=green!22] {angle = A--B--C};
        \node[anchor=north] at (-0.5,5) [font=\small] {$\alpha=\ang{45}$};
        \draw[red, thick] (-1.5, 2.3)--(0, 1.7);
        \draw[red, thick] (0,1.7)--(2.5, 1.2);
        \draw[red, thick] (2.5, 1.2)--(4, 0.6);
        \draw[dashed] (-1.5, 1.7)--(0,1.7);
        \draw[dashed] (2.5, 1.2)--(4,1.2);
        \node[anchor=north] at (-0.5,5) [font=\small] {$\alpha=\ang{45}$};
        \pic [draw, angle radius=1cm, "$\theta$", angle eccentricity=1.6, fill=red!22] {angle = G--F--E};
        \pic [draw, angle radius=1cm, "$\theta$", angle eccentricity=1.6, fill=red!22] {angle = J--I--H};
        % --- AGGIUNTA: Prolungamento raggio incidente per evidenziare d ---
        \draw[dotted, black, thick] (0,1.7) -- (4, 0.1);
        
        % --- AGGIUNTA: Evidenziazione distanza d tra le parallele ---
        \draw[<->, blue] (4, 0.6) -- (3.827586207, 0.1689655172);
        \node[blue, right] at (4.03, 0.3) {$d$};
        
        % --- AGGIUNTA: Spessore h sul lato più in basso ---
        \draw[<->, >=stealth] (0, -0.3) -- (2.5, -0.3);
        \node[below] at (1.25, -0.3) {$h$};
        
    \end{tikzpicture}
    \caption{Schema della deviazione di un raggio incidente in un prisma a base trapezoidale}
\end{figure}

Si osserva che il raggio uscente dal prisma risulta parallelo al raggio incidente, spostato di una distanza d che dipende dall'angolo di incidenza secondo la relazione
\begin{equation}
    \label{eq:spostamento-facce}
    d = h \sin\theta \left( 1 - \frac{\cos\theta}{\sqrt{n^2 - \operatorname\sin^2\theta}} \right)
\end{equation}
Sono state raccolte le misure dello spostamento per diversi angoli di incidenza. Ciascuna è stata confrontata con il valore teorico previsto dalla relazione \eqref{eq:spostamento-facce}, con spessore del prisma trapezoidale $h=(32\pm1)$ mm e indice di rifrazione determinato nelle misure dell'angolo di deviazione minima.
Si raccolgono i dati e le compatibilità nella Tabella 4.

\begin{table}[H]
\centering
\begin{tabular}{cccc}
\toprule
\textbf{Angolo \small{(\unit{\degree})}} & \textbf{Distanza dal centro \small{(\unit{mm})}} & \textbf{Distanza attesa \small{(\unit{mm})}} & $\mathbf{t_0}$ \\
\midrule
10 & \qty{2 \pm 1}{}  & \qty{1.97 \pm 0.22}{}  & 0.03 \\
20 & \qty{4 \pm 1}{}  & \qty{4.08 \pm 0.28}{}  & 0.07 \\
30 & \qty{6 \pm 1}{}  & \qty{6.46 \pm 0.36}{}  & 0.43 \\
40 & \qty{9 \pm 1}{}  & \qty{9.27 \pm 0.47}{}  & 0.25 \\
50 & \qty{12 \pm 1}{} & \qty{12.68 \pm 0.60}{} & 0.58 \\
\bottomrule
\end{tabular}
\caption{Confronto tra le distanze dal centro misurate sperimentalmente e i valori teorici attesi al variare dell'angolo d'incidenza.}
\label{tab:distanze_angoli}
\end{table}

Le incertezze sulle misure sono dovute alla sensibilità della carta millimetrata e del goniometro posti sul banco ottico. 
L'errore sulla distanza attesa è stato ottenuto propagando le incertezze di \ang{1} sulla misura degli angoli di incidenza, di \qty{1}{\milli\meter} sulla larghezza $h$ del prisma e quella dell'indice di rifrazione.
Si osserva una piena compatibilità delle misure entro una deviazione standard.




